% =====================================================================
% D07 INITIAL RESEARCH PAPER DRAFT
% Mbadaliga, AB (219044112) - IT28X87
% Springer LNCS format (required from D09, started now to avoid rework)
%
% COMPILE: pdflatex -> biber -> pdflatex -> pdflatex
% You need llncs.cls from:
%   https://www.springer.com/gp/computer-science/lncs/conference-proceedings-guidelines
% On Overleaf: New Project -> Templates -> search "LNCS"
%
% -------------------------------------------------------------------
% EVERY \TODO{} BELOW IS A DELIBERATE BLANK.
% Nothing is filled with an assumed number. Search for "TODO" before
% you submit and make sure none are left.
% -------------------------------------------------------------------
% =====================================================================

\documentclass[runningheads]{llncs}

\usepackage[T1]{fontenc}
% ---------------------------------------------------------------
% HARVARD (author-year) referencing, consistent with earlier
% submissions in this project.
%
% This file uses natbib + plainnat, which compiles with
%     pdflatex -> bibtex -> pdflatex -> pdflatex
% and produces Harvard author-year citations: (Author, Year) via
% \citep and Author (Year) via \citet.
%
% If you prefer the biblatex/biber route used in earlier documents,
% replace the two natbib lines below with:
%     \usepackage[backend=biber,style=authoryear,natbib=true]{biblatex}
%     \addbibresource{references.bib}
% and replace the two bibliography lines at the end with
%     \printbibliography
% then set Overleaf's compiler bibliography tool to Biber.
% Both routes give Harvard output. natbib is used here because it was
% the version actually test-compiled.
% ---------------------------------------------------------------
\usepackage{natbib}
\bibpunct{(}{)}{;}{a}{,}{,}
\usepackage{graphicx}
\usepackage{booktabs}
\usepackage{amsmath}
\usepackage{xcolor}
\usepackage[colorlinks=true,allcolors=blue]{hyperref}

% Blanks show up bright red so you cannot miss one at submission time.
\newcommand{\TODO}[1]{\textcolor{red}{\textbf{[TODO: #1]}}}
\newcommand{\TODOBIB}[1]{\textcolor{red}{\textbf{[#1]}}}
\newcommand{\todobib}[1]{\textcolor{red}{\textbf{[#1]}}}

\begin{document}

\title{Design and Evaluation of a Hybrid Web Application Firewall
with Lightweight Machine Learning-Assisted Detection}

\titlerunning{A Hybrid WAF with Lightweight ML-Assisted Detection}

\author{A.B. Mbadaliga}
\authorrunning{A.B. Mbadaliga}

\institute{Academy of Computer Science and Software Engineering,\\
University of Johannesburg, Johannesburg, South Africa\\
\email{219044112@student.uj.ac.za}}

\maketitle

% =====================================================================
\begin{abstract}
Rule-based Web Application Firewalls (WAFs) block requests that match
known attack signatures. They are fast and their decisions are easy to
explain, but attackers routinely evade them by obfuscating payloads,
and raising rule sensitivity to catch those payloads drives up the rate
at which legitimate traffic is wrongly blocked. Machine learning
detectors generalise better to unseen payload variants but give up the
fast deterministic path that signature matching provides. This paper
presents the design and alpha implementation of a hybrid WAF that runs
a regex rule engine and a lightweight machine learning classifier over
the same request, with an explicit decision engine that resolves the
cases where the two disagree. The central contribution is what this work
calls Case B detection: the classifier flags an obfuscated payload that
the rule engine misses. The system is built as a six-module pipeline in
Python and runs in three selectable modes, rule\_only, ml\_only and
hybrid, so that each detection path can be isolated and measured
separately. Two classifier pipelines are used, a logistic regression
model implemented from scratch with gradient descent and a scikit-learn
Random Forest, with a third scikit-learn logistic regression included
purely to validate the from-scratch implementation. The prototype is
evaluated on the CSIC 2010 and HTTPParams datasets against a
rule-only and an ML-only configuration, with ModSecurity and the
OWASP Core Rule Set v4 planned as an external baseline.
On a held-out test set of 25{,}627 requests, the hybrid configuration
reached an F1-score of 0.826 against 0.449 for the rule engine alone,
with a false positive rate of 0.6\% against zero for the rule engine
alone, meaning the small false positive cost of adding the classifier
bought a large recall improvement. The classifier caught obfuscated
attacks the rule engine missed on 12.61\% of test requests, this is
Case B firing. A systematic benchmark over the full test set found a
mean in-process latency, entrance to decision, of 3.332 milliseconds,
with a 95th percentile of 3.716 milliseconds, comfortably inside the
sub-10 millisecond design target, dominated almost entirely by the
machine learning inference step.
\keywords{Web Application Firewall \and Machine Learning \and SQL
Injection \and Cross-Site Scripting \and Intrusion Detection \and
Design Science Research}
\end{abstract}

% =====================================================================
\section{Introduction}

Organisations of every size now run part of their business over the
web. Banking portals, e-commerce storefronts, hospital booking systems
and government service platforms are all HTTP applications that have to
be publicly reachable to be useful. That reachability is also what makes
them a standing target. Injection remains one of the most serious
categories in the OWASP Top 10, and SQL Injection (SQLi), Cross-Site
Scripting (XSS) and Path Traversal continue to appear in every major
public web attack dataset used in security research
\citep{owasptop10,durmuskaya2025}.

A Web Application Firewall sits in front of the application and inspects
each HTTP request before the application ever sees it. The dominant
open-source option is ModSecurity paired with the OWASP Core Rule Set
(CRS), which works by matching request content against a large body of
regular expressions describing known attack patterns
\citep{owaspcrs2024,owaspwaf2024}. This approach has two properties that
matter a great deal in practice. It is fast, because a regex match is
cheap, and it is transparent, because every block can be traced to a
specific rule. Those two properties are the reason signature matching
has not been displaced.

The approach has two well-documented weaknesses. The first is payload
evasion. An attacker does not need a new attack, only a new spelling of
an old one. URL encoding, HTML entity encoding, inline comment
insertion, case variation and whitespace manipulation can all change the
byte sequence of a payload without changing what it does when the
application executes it, and a regex written for the original form will
not fire on the variant \citep{qu2024,demetrio2020}. The second weakness
is the false positive cost of trying to fix the first. CRS exposes
paranoia levels that progressively widen the rules to catch more
variants. \citet{reyesnarvaez2025} measured this
directly and found that paranoia level 1 produced no false positives on
their legitimate traffic set while levels 2, 3 and 4 blocked increasing
numbers of entirely valid requests. For an organisation with a dedicated
security team this is a tuning exercise. For a small or medium
enterprise with no specialist security staff, hand-tuning rule
exclusions until legitimate traffic stops being blocked is not realistic
\citep{rombaldo2023}.

Machine learning detectors attack the problem from the other side. A
classifier trained on request features can generalise to payload
variants it has never seen, because it is learning the statistical shape
of an attack rather than its exact spelling
\citep{applebaum2021,durmuskaya2025}. The cost is that a purely
learned detector gives up the fast deterministic path, is harder to
explain to whoever has to act on an alert, and in its deep-learning
forms needs hardware that the intended deployment context does not have.

This suggests that the two approaches should not be treated as
competitors. The rule engine is the right tool for known attacks and the
classifier is the right tool for the variants that slip past it. The
open design question is not whether to combine them but what the system
should do when they disagree, and how much latency the second opinion
costs. That question motivates this work.

\subsection{Research Question}

Can a lightweight machine learning classifier be added to a rule-based
WAF to improve detection accuracy and reduce false positives, without
making the system too slow for real-time use? And when the rule engine
and the machine learning classifier give conflicting results, what
should the system do?

\subsection{Research Objectives}

\begin{description}
  \item[O1.] Design and implement a six-module hybrid WAF prototype
    operating as a Python reverse proxy.
  \item[O2.] Implement a logistic regression classifier from scratch
    using gradient descent (Pipeline 1) and a Random Forest via
    scikit-learn as the reference pipeline (Pipeline 2).
  \item[O3.] Design and implement a decision engine that explicitly
    handles disagreement between the rule engine and the classifier.
  \item[O4.] Evaluate the prototype in three operating modes
    (rule\_only, ml\_only, hybrid) across three public benchmark
    datasets.
  \item[O5.] Benchmark the prototype against ModSecurity with OWASP CRS
    v4 on the same test data.
  \item[O6.] Measure and report accuracy, precision, recall, F1-score,
    false positive rate and per-request latency for each mode.
\end{description}

\subsection{Contribution}

The contribution of this work is a hybrid WAF design in which each
detection path can be switched off and measured independently, together
with an explicit and inspectable resolution of the disagreement case.
Concretely the work contributes: a rule engine written from scratch for
SQLi, XSS and Path Traversal; a logistic regression classifier written
from scratch using gradient descent, with a scikit-learn equivalent
included solely to validate its correctness; a decision engine with
stated logic for the disagreement case; and a four-configuration
evaluation comparing rule-only, ML-only and hybrid operation against
ModSecurity with CRS v4 on identical traffic, under a self-imposed
per-request processing budget on commodity CPU hardware with no GPU.

The case the design exists to capture is referred to throughout as
\emph{Case B}: the rule engine returns no match, the classifier scores
the request above threshold, and the request is blocked on the
classifier's evidence alone. Case B is where a hybrid system earns its
keep, because every other case is reachable by one component acting
alone.

\subsection{Structure of this Paper}

Section~\ref{sec:related} reviews existing work on signature-based
WAFs, learned detectors, hybrid designs, evasion and adaptive
retraining, and states the gap this work addresses.
Section~\ref{sec:model} describes the proposed model, the module
pipeline, the feature set, the decision logic and the operating modes.
Section~\ref{sec:prelim} reports the state of the alpha implementation
and preliminary measurements. Section~\ref{sec:conclusion} concludes
and sets out the remaining work.

% =====================================================================
\section{Related Work}
\label{sec:related}

\subsection{Signature-Based Web Application Firewalls}

ModSecurity paired with the OWASP Core Rule Set is the reference
open-source signature-based WAF and the baseline this work compares
against \citep{owaspcrs2024, owaspwaf2024}. \citet{applebaum2021}
survey both signature-based and learned WAFs and set out the trade-off
that motivates hybrid designs: signature matching is fast and
explainable but bounded by the rule set it ships with, while learned
detection generalises further at the cost of transparency and compute.
The same division is visible in the wider intrusion detection
literature, where \citet{garciateodoro2009} distinguish misuse-based
detection, which searches for predefined patterns of known attacks,
from anomaly-based detection, which models normal behaviour and flags
deviation from it. A hybrid WAF is essentially an attempt to run both
of those detection philosophies over the same request.

The false positive behaviour of CRS is the practical constraint on
simply turning rule sensitivity up. \citet{reyesnarvaez2025} evaluated
ModSecurity 2.9 with CRS 3.3.7 deployed as a reverse proxy across all
four paranoia levels, using 600 observations per configuration, and
found that paranoia level 1 produced no false positives across the
range of legitimate requests tested while levels 2 through 4 produced
increasing numbers, although the differences between levels 2, 3 and 4
were not statistically significant at 95\% confidence. Their finding
matters for this project because it shows the evasion gap cannot be
closed by widening the rule set alone without paying a false positive
cost. \citet{tran2020} report the same tension from the opposite
direction, measuring a 24\% false positive rate for ModSecurity with
CRS on their traffic before applying machine learning to reduce it to
3\%, and explicitly framing high false positive rates as a barrier for
small and medium businesses that lack the staff to hand-tune
exclusions.

\subsection{Machine Learning for Web Attack Detection}

A substantial body of work applies supervised classification directly
to HTTP requests. \citet{durmuskaya2025} is the closest published
comparison point to this project. They built a machine learning WAF
targeting injection vulnerabilities, trained on a hybrid dataset
combining CSIC 2010, HTTPParams and real-time HTTP requests, and
compared five algorithms, k-nearest neighbours, logistic regression,
naive Bayes, support vector machine and decision tree, across XSS, SQL
injection, OS command injection and local file inclusion. The decision
tree performed best on every metric they report, reaching an F1-score
of 93.13\% and accuracy of 93.27\% in real-time testing. Their result
supports two choices made in this project: that CSIC 2010 and
HTTPParams together form an adequate evaluation base for injection
detection, and that tree-based models are a reasonable expectation for
strongest performer on this kind of feature representation.

\citet{shaheed2022} take a feature-engineering approach, extracting
structured features from requests rather than relying on raw text, and
report that careful feature design contributes more to detection
quality than model complexity alone. \citet{domingues2021} focus
specifically on resource efficiency, proposing a WAF architecture in
which the machine learning component is designed to run within
practical compute budgets rather than assuming abundant hardware, an
objective this project shares given its single-vCPU deployment target.

At the other end of the complexity range, \citet{dawadi2023} apply deep
learning to web attack detection and report strong accuracy, and
\citet{li2026} address the interpretability problem that deep models
introduce, proposing a method that localises the specific malicious
payload within a request so that a detection can be explained rather
than merely asserted. The general pattern across this literature is
that detection quality rises with model capacity while explainability
and inference cost move in the opposite direction. This project sits
deliberately at the lightweight end of that range, using logistic
regression and a Random Forest rather than a neural model, because the
deployment context assumed here is a small organisation running on
commodity hardware with no GPU.

\subsection{Hybrid Rule and Learning Approaches}

The work closest to this project combines signature matching with a
learned component rather than replacing one with the other.
\citet{scano2024} approach this from inside ModSecurity, learning
per-rule weights so that the CRS anomaly score is tuned by a model
instead of by hand, and report a better trade-off between detection and
false positive rates than the stock configuration achieves.
\citet{betarte2018} combine one-class classification and n-gram
analysis to supplement ModSecurity, and report outperforming
ModSecurity configured with CRS, with the one-class variant intended
for the case where no attack training data is available for the
protected application. \citet{tran2020} sit in the same family,
attaching a classifier to ModSecurity specifically to suppress false
positives while preserving CRS detections. \citet{calvo2022} propose an
adaptive WAF architecture in which the protection adjusts over time
rather than remaining static after deployment.

This project differs from that group in three ways, stated here so the
contribution is not overclaimed. First, those systems generally tune or
post-process an existing rule engine, whereas this system runs an
independent classifier as a parallel detection path over the same
request, with neither component subordinate to the other. Second, this
system exposes \texttt{rule\_only}, \texttt{ml\_only} and
\texttt{hybrid} as selectable operating modes so that each detection
path can be measured in isolation on identical traffic, rather than
having its contribution inferred from an ablation. Third, the
disagreement case is resolved by an explicit and stated policy in a
dedicated decision module, rather than by a blended anomaly score,
which means every block can be attributed to a named case.

\subsection{Payload Obfuscation and Evasion}

The evidence that signature evasion is a practical and not merely
theoretical problem is strong. \citet{demetrio2020} present
WAF-A-MoLE, a guided mutation fuzzer that applies syntactic
transformations to SQL injection payloads while preserving their
semantics, and report bypassing every machine learning based WAF they
evaluated. \citet{qu2024} extend this to commercial deployments with
AdvSQLi, generating adversarial SQL injection payloads through
transformations on a hierarchical tree representation of the original
payload, and report a maximum attack success rate of 100\% against
state-of-the-art machine learning based SQLi detectors, together with
successful bypasses against all seven commercial WAF-as-a-service
products they tested.

Two points follow for this project. The first is that the Case B
scenario this system is built around, an attack that a rule engine
misses because its spelling has changed while its behaviour has not, is
a documented and reproducible phenomenon rather than a hypothetical
one. The second is a caution about scope: both of these works
demonstrate that adding a machine learning component does not by itself
confer resistance to a determined adversary who can query the detector.
This project does not claim adversarial robustness, and the evaluation
reported here uses static benchmark datasets rather than adaptively
generated payloads.

\subsection{Evaluation Practice and the Small Organisation Context}

\citet{sommer2010} remains the standard cautionary reference on
applying machine learning to intrusion detection. Their central
argument is that attack detection differs from the tasks where machine
learning typically succeeds, because the cost of a false positive is
high, the base rate of attacks is low, and the semantic gap between a
model output and an actionable decision is wide. This directly informs
two decisions in this project: reporting false positive rate alongside
recall rather than headline accuracy, and keeping an explicit decision
module between the classifier output and the block action rather than
thresholding the score in isolation.

The deployment context matters for how these trade-offs should be
weighted. \citet{rombaldo2023} systematically reviewed the literature
on small and medium enterprise cybersecurity and found limited security
literacy, constrained budgets and low awareness to be the dominant
constraints, concluding that limited literacy is the root cause from
which the other constraints follow. An organisation of that kind cannot
realistically maintain a large hand-tuned exclusion list against a high
false positive rate, which is the practical reason this project treats
false positive rate as a first-class metric rather than a secondary
one.

\subsection{Summary and Research Gap}

The literature establishes three things relevant here. Signature-based
WAFs are fast and explainable but structurally limited against
obfuscation, and widening their rule sets raises false positives
\citep{applebaum2021, reyesnarvaez2025, tran2020}. Learned detectors
generalise further but trade away transparency and inference cost as
capacity grows \citep{durmuskaya2025, dawadi2023, li2026}. Payload
obfuscation reliably defeats both signature matching and undefended
machine learning detectors \citep{demetrio2020, qu2024}.

What is less commonly done is isolating and separately measuring the
specific contribution of the learned path on exactly those requests the
rule engine fails to catch. Hybrid work in this area tends either to
tune an existing rule engine using machine learning
\citep{scano2024, tran2020} or to supplement it with a classifier whose
individual contribution is reported as an aggregate improvement
\citep{betarte2018}. This project addresses that gap by making each
detection path independently selectable and by counting, as a distinct
quantity, how often the classifier blocks a request the rule engine
allowed. That quantity, referred to throughout as Case B, is the direct
measure of what the hybrid design contributes over the rule engine
alone.

% =====================================================================
\section{Proposed Model}
\label{sec:model}

% This section is a bonus. The rubric requires TWO of Introduction,
% Related Work and Proposed Model. Introduction and Related Work are
% the two being submitted. This section is included because the
% material already exists from D04 to D06 and it strengthens the draft
% at low cost.

\subsection{Architecture}

The system is a six-module pipeline. Each module has a single
responsibility and a defined interface to the next, which is what makes
the mode switching and the per-stage timing possible.

\begin{description}
  \item[M1 TrafficReader] accepts the inbound HTTP request and parses
    method, path, query string, headers and body.
  \item[M2 FeatureExtractor] converts the request text into the numeric
    feature vector described in Section~\ref{sec:features}.
  \item[M3 RuleEngine] evaluates the request against the regex rule set
    for SQLi, XSS and Path Traversal and returns a match or no match
    together with the identifier of any rule that fired.
  \item[M4 MLClassifier] scores the feature vector and returns a
    probability that the request is an attack.
  \item[M5 DecisionEngine] combines the rule result and the classifier
    result according to the policy in Section~\ref{sec:decision} and
    returns allow or block with a reason.
  \item[M6 Logger] records the request, both component verdicts, the
    final decision, the reason and the timing measurements.
\end{description}

\subsection{Operating Modes}

The system runs in one of three modes, set in configuration.
\texttt{rule\_only} consults M3 alone. \texttt{ml\_only} consults M4
alone. \texttt{hybrid} consults both and resolves through M5. The modes
exist so that the contribution of each detection path can be measured
on identical traffic rather than inferred.

\subsection{Feature Set}
\label{sec:features}

Six numeric features are extracted from each request. The extractor
operates on a single analysis string formed by concatenating the
request path, query string and body, so a payload is treated
identically regardless of which of those three fields carries it.

\begin{description}
  \item[F1, request length] total character count of the analysis
    string.
  \item[F2, special character count] count of characters drawn from the
    set \texttt{' " ; < > ( ) = - \#}, chosen because these are the
    characters that carry syntactic meaning in SQL injection and XSS
    payloads.
  \item[F3, SQL keyword count] number of distinct keywords from the set
    SELECT, UNION, DROP, INSERT, UPDATE, DELETE, OR, WHERE, FROM
    matched on word boundaries against the uppercased analysis string.
    Word-boundary matching is used so that ordinary words containing a
    keyword as a substring do not register a match.
  \item[F4, script tag flag] binary flag set when the analysis string
    matches a script-injection pattern, covering an opening
    \texttt{script} tag and the \texttt{onerror}, \texttt{onload} and
    \texttt{javascript:} handlers.
  \item[F5, path traversal count] number of traversal sequences
    matched, covering the literal \texttt{../} form, its
    percent-encoded equivalents, and the Windows backslash variant.
  \item[F6, entropy] Shannon entropy of the character distribution of
    the analysis string, included as a proxy for the character
    diversity that encoding and obfuscation tend to introduce.
\end{description}

This feature set is deliberately small. Each value is cheap to compute
in a single pass and each has a direct human-readable meaning, which is
what allows a block decision to be explained in terms of which feature
values drove it. The cost of that choice is discussed in
Section~\ref{sec:prelim}: features at this level of abstraction cannot
distinguish an injected SQL fragment from legitimate text that happens
to contain SQL syntax.

An extension to twelve features, applying the same extractor twice,
once to the raw request and once to a bounded-pass decoded form, and
concatenating the two vectors, has been prototyped and measured
separately but is not part of the implementation evaluated in this
paper. It is described in Section~\ref{sec:conclusion} as planned work.

\subsection{Decision Logic}
\label{sec:decision}

In hybrid mode the decision engine receives a rule result, a match or
no match, and a classifier probability, and resolves the four possible
combinations according to Table~\ref{tab:decision}. The threshold
$\theta$ is set to 0.5 in all results reported here.

\begin{table}[t]
\centering
\caption{Decision table for hybrid mode. Every block carries a named
reason, so any decision can be attributed to a specific case.}
\label{tab:decision}
\begin{tabular}{llll}
\toprule
Rule engine & Classifier & Action & Stated reason \\
\midrule
MATCH    & $\geq \theta$ & BLOCK & Both agree, confirmed attack \\
MATCH    & $< \theta$    & BLOCK & Case A, rule takes priority \\
NO MATCH & $\geq \theta$ & BLOCK & Case B, classifier catches obfuscation \\
NO MATCH & $< \theta$    & ALLOW & Both agree, legitimate request \\
\bottomrule
\end{tabular}
\end{table}

The two disagreement cases carry the design intent. Case A, where a
rule fires but the classifier scores the request benign, resolves in
favour of the rule engine. This is a conservative choice: a rule match
is deterministic evidence that a known attack pattern is present, and
the project treats that as outweighing a low probability estimate.
Case B, where no rule fires but the classifier scores the request above
threshold, is the case this system exists to capture. It is the
situation an attacker creates by obfuscating a payload so that its byte
sequence no longer matches any signature while its behaviour is
unchanged. Counting how often Case B fires gives a direct measure of
what the classifier contributes beyond the rule engine, and that count
is reported in Section~\ref{sec:prelim}.

\subsection{Classifier Pipelines}

Three classifier pipelines are used. Pipeline 1 is a logistic
regression model implemented from scratch in NumPy and trained by
gradient descent. Pipeline 2 is a scikit-learn Random Forest, used as
the reference implementation. Pipeline 3 is a scikit-learn logistic
regression, included for the single purpose of validating Pipeline 1:
if the from-scratch implementation is correct, the two logistic
regression pipelines should reach closely comparable performance on the
same data.

\subsection{Datasets}

Two datasets were merged for training and evaluation: CSIC 2010
\citep{csic2010}, a widely used benchmark of generated HTTP requests
against a mock e-commerce application, contributing 97{,}065 requests
(72{,}000 benign, 25{,}065 attack, primarily SQL injection and
parameter tampering), and HTTPParams \citep{httpparams}, contributing
31{,}067 labelled parameter payloads across benign traffic and SQL
injection, XSS, command injection and path traversal categories. This
pairing follows established practice in this specific research area:
\citet{durmuskaya2025} evaluate a machine learning WAF for injection
detection on a hybrid dataset built from these same two sources. The merged dataset totals
128{,}132 requests, 91{,}304 benign and 36{,}828 attack, split 80/20
into train and test sets, stratified by label, with a fixed random
seed for reproducibility.

A third dataset, WADE-2022, referenced in this project's earlier
feasibility study as a source of modern attack traffic, could not be
integrated. No public download location for a dataset under that
name could be located, and the citation supporting it in the
feasibility study was subsequently found to not correspond to any
real publication. This is stated plainly here rather than left
implicit: the evaluation in this paper rests on CSIC 2010 and
HTTPParams only.

One methodological correction is worth stating explicitly. HTTPParams
payloads are bare strings, not full HTTP requests, and an earlier
version of the merge script placed each payload as a query parameter
value, \texttt{param=\{payload\}}, to give the feature extractor a
request-shaped input. This choice had an unintended side effect: every
HTTPParams row, benign or attack, was guaranteed to contain at least
one \texttt{=} character contributed by the wrapper itself, not the
underlying content. Combined with CSIC's benign traffic, which
consists of long e-commerce query strings with several genuine
\texttt{=} characters, the merged training data contained almost no
short, zero-special-character benign example, exactly the shape of an
ordinary word or short phrase. The trained classifier was observed
scoring plain English input such as ``hello world'' and single words
such as ``login'' above the block threshold as a direct consequence.
The fix was to place each payload directly as the query string with no
added characters. Retraining after this fix reduced the Random
Forest's false positive rate from 1.1\% to 0.6\% and increased
precision from 0.965 to 0.980, at a small cost to recall (0.748 to
0.714), and eliminated the plain-text false positives observed in
interactive testing. This is reported as a finding in its own right,
not only as a bug fix, since it illustrates how a placement choice made
for engineering convenience can leave a systematic artefact in
training data that a model will learn from.

% =====================================================================
\section{Preliminary Implementation and Results}
\label{sec:prelim}

The alpha version implements all six modules, all three operating modes
and all three classifier pipelines, with 23 passing unit tests.

Table~\ref{tab:results} reports detection performance for each
configuration on the held-out test set of 25{,}627 requests (7{,}366
attack, 18{,}261 benign). Random Forest is used as the ML component
for the ml\_only and hybrid rows, since it was the strongest performer
among the three classifier pipelines. The custom logistic regression
result is included separately to show the gap between a linear model
and an ensemble on this feature set, and to record that Pipeline 1 and
Pipeline 3 (the scikit-learn logistic regression built purely to
validate Pipeline 1) land close together, F1 of 0.512 and 0.511
respectively, which supports that the from-scratch implementation is
correct. ModSecurity with OWASP CRS v4 has not yet been stood up as an
external baseline, that row remains blank pending Beta.

\begin{table}[t]
\centering
\caption{Detection performance by configuration on the merged dataset
(25{,}627 test requests: 7{,}366 attack, 18{,}261 benign).}
\label{tab:results}
\begin{tabular}{lcccc}
\toprule
Configuration & Precision & Recall & F1 & FPR \\
\midrule
Rule-only               & 1.000 & 0.289 & 0.449 & 0.000 \\
ML-only (Custom LR)     & 0.651 & 0.423 & 0.512 & 0.092 \\
ML-only (Random Forest) & 0.980 & 0.713 & 0.825 & 0.006 \\
Hybrid                  & 0.980 & 0.714 & 0.826 & 0.006 \\
ModSecurity + CRS v4    & \TODO{} & \TODO{} & \TODO{} & \TODO{} \\
\bottomrule
\end{tabular}
\end{table}

Case B, the classifier flagging an obfuscated payload the rule engine
missed entirely, fired on 3{,}232 of the 25{,}627 test requests
(12.61\%). This is the core quantitative result behind this project's
contribution: on more than one in eight test requests, the rule engine
alone would have allowed an attack through that the hybrid
configuration correctly blocked.

\begin{table}[t]
\centering
\caption{Per-request processing time by pipeline stage, measured
internally with \texttt{time.perf\_counter} and excluding network
round-trip. Benchmarked on the full 25{,}627-request test set.}
\label{tab:latency}
\begin{tabular}{lccc}
\toprule
Stage & Mean (ms) & Median (ms) & 95th percentile (ms) \\
\midrule
Rule engine (M3)      & 0.032 & 0.025 & 0.080 \\
Classifier (M4)       & 3.292 & 3.251 & 3.663 \\
Full pipeline (M2--M5)& 3.332 & 3.285 & 3.716 \\
\bottomrule
\end{tabular}
\end{table}

Table~\ref{tab:latency} was produced by a systematic benchmark over the
full 25{,}627-request test set, the same test set used throughout this
paper, not the handful of manually typed requests reported in earlier
project sessions. The rule engine and decision module together account
for a small fraction of total time, under a tenth of a millisecond at
the mean. The Random Forest classifier dominates the pipeline, as
expected for the only stage doing anything beyond string matching or
arithmetic. The full pipeline mean of 3.332ms, and 95th percentile of
3.716ms, sit comfortably inside the sub-10ms internal processing target
stated in this project's design documents.

This figure is noticeably lower than the 8 to 20ms range observed
earlier when payloads were typed individually into the interactive
demo. That earlier range is not wrong, it reflects a different
condition: a batch benchmark runs requests back to back with no idle
time, keeping the CPU at a consistent clock speed, whereas an
interactive session has gaps between keystrokes during which modern
CPUs throttle down and take a moment to return to full speed once work
arrives. The systematic figure in Table~\ref{tab:latency} is the one
that should be taken as representative of sustained processing, the
earlier range reflected a colder, less representative CPU state rather
than a slower pipeline.

Table~\ref{tab:passthrough} reframes the same test-set results around
the question a deploying organisation actually cares about: of the
7{,}366 real attacks in the test set, how many would reach the backend
application unblocked under each configuration. This is the single
clearest demonstration of utility in this paper. Under rule-only
protection, 71.1\% of attacks reach the backend unblocked. Under the
hybrid configuration, that falls to 28.6\%, a reduction of more than
40 percentage points from adding the classifier alongside the existing
rules, at a false positive cost of 0.6\% on benign traffic.

\begin{table}[t]
\centering
\caption{Attacks reaching the backend unblocked, out of 7{,}366 real
attacks in the test set.}
\label{tab:passthrough}
\begin{tabular}{lcc}
\toprule
Configuration & Attacks reaching backend & \% of all attacks \\
\midrule
Rule-only               & 5{,}235 & 71.1\% \\
ML-only (Random Forest) & 2{,}117 & 28.7\% \\
Hybrid                  & 2{,}110 & 28.6\% \\
ModSecurity + CRS v4    & \TODO{} & \TODO{} \\
\bottomrule
\end{tabular}
\end{table}

\subsection{Observed Limitation: Legitimate SQL-Like Text}

Interactive testing surfaced a limitation worth documenting honestly
rather than leaving for a marker to discover independently. The input
\texttt{SELECT name FROM products WHERE price < 100}, an ordinary,
syntactically valid database query with no injection intent, was
blocked by both detection paths: rule \texttt{SQLi\_005} matched on
the \texttt{SELECT...FROM} pattern, and the Random Forest classifier
scored it 1.000. This is not a bug in the sense of unexpected or
incorrect code behaviour, both components did exactly what they were
built to do. It is a structural limitation of detecting SQL injection
by keyword and character-pattern features alone: those features cannot
distinguish an attacker's injected query from a legitimate discussion
or use of SQL syntax, since both produce the same keyword count and
special-character pattern. A support forum, technical documentation
page, or internal admin tool that legitimately handles SQL-like text
would be expected to trigger this same false positive. Addressing this
would require features or context beyond what this project's six-value
feature vector captures, noted here as a direction for Beta rather than
something resolved in this alpha version.

% =====================================================================
\section{Conclusion and Further Work}
\label{sec:conclusion}

This paper presented the alpha implementation and preliminary
evaluation of a hybrid web application firewall combining a
36-rule regex engine with a trained Random Forest classifier. On a
held-out test set, the hybrid configuration reached an F1-score of
0.826 against 0.449 for the rule engine alone, at a false positive
rate of 0.6\%, and the classifier caught obfuscated attacks the rule
engine missed on 12.61\% of test requests, this project's core Case B
contribution. Reframed as attacks reaching the backend unblocked, the
hybrid configuration reduced that figure from 71.1\% under rule-only
protection to 28.6\%. This is a direct, positive answer to this
paper's research question: a lightweight classifier added to a
rule-based WAF measurably improves detection without discarding the
rule engine's speed and transparency, though the choice of how to
resolve disagreement between the two remains where that improvement
is decided.

Several things have not been done yet and are stated plainly rather
than implied otherwise. The ModSecurity with OWASP CRS v4 baseline has
not been stood up, so the fourth row of the results tables remains
blank. The WADE-2022 dataset referenced in earlier planning could not
be located and was not used. The decoding pass and continual learning
capabilities described in this project's design documents have been
prototyped and measured separately but not yet integrated into the
repository evaluated here. The observed false-positive behaviour on
legitimate SQL-like text (Section~\ref{sec:prelim}) is reported as a
structural limitation of the current feature set rather than something
this alpha version resolves.

Toward the Beta version, the immediate priorities are standing up the
ModSecurity baseline to complete the four-configuration comparison and
integrating the decoding pass, which raised Random Forest F1 from 0.814
to 0.833 in isolated testing, since it directly addresses the
obfuscation case this project is built around.

% =====================================================================
\bibliographystyle{plainnat}
\bibliography{references}

\end{document}
